% joint-vs-conditional.tex — Joint vs conditional likelihoods for TKF EM
% To be \input from lhopital-limits.tex

\subsection{Joint vs Conditional Likelihoods}
\label{sec:joint-vs-conditional}

The TKF91 transition matrix includes factors of $\kappa = \insrate/\delrate$
and $(1{-}\kappa)$ arising from the geometric$(\kappa)$ prior on
ancestor length: $\Pr(|\text{anc}| = i) = \kappa^i(1{-}\kappa)$.
This prior is part of the \emph{joint} likelihood $P(x,y)$
but can be factored out with respect to the ancestor-length prior to yield an \emph{ancestor length-conditioned} likelihood
$P_{\mathrm{cond}}(x,y \mid |x|=i)
= P(x,y) / [\kappa^i(1{-}\kappa)]$.

\paragraph{The four regimes.}
Of the four combinations of $\kappa < 1$ vs $\kappa = 1$
and joint vs ancestor length-conditioned, only one is degenerate:
\begin{center}
\begin{tabular}{lcc}
\toprule
& $\kappa < 1$ & $\kappa = 1$ \\
\midrule
Joint $P(x,y)$ & well-defined & \textbf{degenerate} ($\to 0$) \\
Ancestor length-conditioned $P_{\mathrm{cond}}$ & well-defined & well-defined \\
\bottomrule
\end{tabular}
\end{center}
At $\kappa = 1$, the geometric prior places vanishing mass on finite
sequences, so $P(x,y) \to 0$ and
$\log P \to -\infty$.
The ancestor length-conditioned WFST removes the $\kappa$ and $(1{-}\kappa)$
factors corresponding to the stationary length distribution from the transition matrix, leaving only
$(\alpha,\beta,\gamma)$-dependent terms.
These have finite L'H\^opital limits (Sections~\ref{sec:score-limits}--\ref{sec:lhopital-bdi}),
so $P_{\mathrm{cond}}$ is finite and well-defined even at $\kappa = 1$.

\paragraph{BDI sufficient statistics.}
The BDI sufficient statistics  $\mathbb{E}[B]$, $\mathbb{E}[D]$,
$\mathbb{E}[S]$ are most naturally recovered from the ancestor length-conditioned WFST, whose score depends only on the six TKF bridge groups 
$(\alpha, 1{-}\alpha, \beta, 1{-}\beta, \gamma, 1{-}\gamma)$. In the joint model, additional contributions from the ancestral length prior appear through the
$\kappa$ and $1{-}\kappa$ factors. These must be included in the M-step when optimizing the full stationary TKF likelihood, but they are conceptually distinct from the bridge statistics themselves.

Equations~(\ref{eq:exposure-tkf})-\ref{eq:deaths-tkf} give the correct expectations for the BDI sufficient statistics
in terms of the counts $\hat{n}_{ij}$ regardless of whether those counts were computed using the joint Pair HMM or the conditionally-normalized transducer,
as long as the WFST transition matrix $\tkfwfst$ is used for the counts $\emcount^\xi_{ij} = \xi \partial_\xi \log \tkfwfst_{ij}$
(e.g. using the tables in Section~\ref{sec:score-table-general} or Section~\ref{sec:score-table-limits}).
The M-step for the joint likelihood includes additional terms from the $\kappa$ and $1{-}\kappa$ factors, leading to a coupled system of equations that can be solved in closed form (Section~\ref{sec:bw-tkf91}).

\paragraph{Closed-form M-steps.}
Both formulations yield closed-form M-steps:
\begin{itemize}
\item \emph{Ancestor-length conditioned:} pure BDI exponential family.
  For the conditional model away from the $\kappa=1$ singularity of the joint prior,
  setting $\partial Q_{\mathrm{cond}}/\partial\insrate = 0$ gives
  \begin{align}
  \hat\insrate &= \frac{\max(0,\;\mathbb{E}[B] + \alpha_\insrate - 1)}
                       {\mathbb{E}[S] + \evoltime + \beta_\insrate}, &
  \hat\delrate &= \frac{\max(0,\;\mathbb{E}[D] + \alpha_\delrate - 1)}
                       {\mathbb{E}[S] + \beta_\delrate}
  \label{eq:mstep-cond}
  \end{align}
  where $(\alpha,\beta)$ are Gamma prior parameters.

\item \emph{Joint ($\kappa < 1$):}
  includes $L \log\kappa + M \log(1{-}\kappa)$.
  Since, for a single sequence-pair, $L  = i$ (ancestor length, exact) and
  $M  = 1$ (one End transition, exact),
  with corresponding accumulated values across multiple sequence-pairs,
  the optimality conditions
  $\partial Q/\partial\insrate = 0$ and $\partial Q/\partial\delrate = 0$
  form a coupled system in $(\insrate,\delrate)$.
  These can be reduced to a single quadratic equation in $\kappa$ with a closed-form solution, which can be substituted back to get $\hat\insrate$ and $\hat\delrate$,
  as described in Section~\ref{sec:bw-tkf91}.
\end{itemize}
In both cases, the M-step maximizes the correct
$Q$-function, so EM monotonicity is guaranteed for
the corresponding tracked objective
($\log P_{\mathrm{cond}}$ for the ancestor length-conditioned model, $\log P$ for the full joint model).
